\chapter{Yeast Infection (Candidiasis)}

\section*{What Is This Test?}
Yeast infection testing detects \textit{Candida} species, the causative agents of vulvovaginal candidiasis (VVC). \textit{Candida albicans} accounts for 80--92\% of VVC cases; non-albicans species, particularly \textit{C. glabrata} (\textit{Nakaseomyces glabratus}), \textit{C. parapsilosis}, \textit{C. tropicalis}, and \textit{C. krusei} (\textit{Pichia kudriavzevii}), account for 5--20\% and are more common in women with recurrent VVC, diabetes, immunocompromised status, and those who have received multiple courses of azole antifungals. \textit{Candida} species are dimorphic fungi that exist as commensal organisms in the vagina of 20--30\% of healthy asymptomatic women. VVC results from overgrowth of \textit{Candida} (usually endogenous) causing symptomatic infection characterized by vulvovaginal pruritus (the most common symptom), burning, irritation, dysuria (external, when urine contacts the inflamed vulvar epithelium), dyspareunia, and a thick, white, "cottage-cheese"-like discharge that is typically non-malodorous. The vaginal pH is normal (3.8--4.5), which is a key differentiating feature from bacterial vaginosis and trichomoniasis (both elevate the pH above 4.5). Approximately 75\% of women will experience at least one episode of VVC in their lifetime. Recurrent VVC (RVVC), defined as four or more symptomatic episodes in a 12-month period, affects 5--10\% of women.

\section*{Alternative Names}
\begin{itemize}
  \item Vaginal Yeast Test
  \item Candida Vaginitis Testing
  \item KOH Prep
  \item Fungal Culture
  \item Vaginitis Panel
  \item Candida PCR
  \item Yeast Culture and Sensitivity
  \item Wet Mount for Yeast
\end{itemize}
\section*{Principle}
Saline wet mount: A drop of vaginal discharge is mixed with a drop of 0.9\% saline on a glass slide, covered with a coverslip, and examined under 400x magnification for budding yeast (blastoconidia) and pseudohyphae/hyphae (the invasive form associated with symptomatic infection). Sensitivity is only 40--70\%, and a negative wet mount does NOT exclude VVC. KOH (10\%) preparation: A drop of vaginal discharge is mixed with a drop of 10\% KOH, which dissolves epithelial cells and cellular debris, making fungal elements (budding yeast, pseudohyphae, hyphae) more visible. KOH is more sensitive than saline (50--80\%) for detecting yeast. Gram stain: Vaginal smear Gram stained: gram-positive budding yeast and pseudohyphae are visualized. Sensitivity comparable to KOH. Fungal culture: Sabouraud dextrose agar (with gentamicin and chloramphenicol to suppress bacteria) inoculated with vaginal swab and incubated at 30--35\textdegree{}C for 48--72 hours (some protocols hold for up to 5 days). Cream-colored, pasty colonies are identified as \textit{Candida} species. Chromogenic media (CHROMagar Candida) differentiate species by colony color: \textit{C. albicans} (green), \textit{C. tropicalis} (blue), \textit{C. krusei} (pink, rough). Species identification by MALDI-TOF mass spectrometry or biochemical testing. Culture is more sensitive than microscopy (70--90\%) and allows species-level identification and antifungal susceptibility testing if indicated. NAAT: Multiplex PCR panels (BD MAX Vaginal Panel, Hologic Aptima CV/TV, BioFire FilmArray Vaginitis Panel) detect \textit{Candida} species group (\textit{C. albicans}, \textit{C. tropicalis}, \textit{C. parapsilosis}, \textit{C. dubliniensis}) and \textit{C. glabrata} individually with high sensitivity (96--99\%) and specificity (97--99\%). Results in 1--3 hours. NAATs are the most sensitive diagnostic method and are increasingly used in clinical practice.

\section*{History}
Vulvovaginal candidiasis was first described as a clinical entity in the mid-19th century. The KOH preparation was introduced in the early 20th century as a simple, rapid diagnostic method. Nystatin, the first polyene antifungal, was discovered in 1950. The first azole antifungal (miconazole) was introduced in the 1970s, followed by clotrimazole, butoconazole, and others. Oral fluconazole was approved in 1990 and became the preferred single-dose oral therapy for uncomplicated VVC. The emergence of fluconazole-resistant \textit{Candida} species (\textit{C. glabrata}, \textit{C. krusei}, and increasingly \textit{C. albicans}) has driven the need for species-level identification and antifungal susceptibility testing in recurrent and refractory cases. Multiplex molecular panels for vaginitis were introduced in the 2010s, incorporating \textit{Candida} detection alongside BV and \textit{T. vaginalis} assays.

\section*{Why This Test Is Ordered}
\begin{itemize}
  \item Diagnosis of acute VVC in a woman with vulvovaginal pruritus, burning, irritation, dysuria, and thick white "cottage-cheese"-like discharge
  \item Differentiation of VVC from other causes of vaginitis: bacterial vaginosis (foul odor, homogeneous gray discharge, pH >4.5, clue cells), trichomoniasis (frothy yellow-green discharge, pH >4.5, motile trichomonads), and mixed infections (which are common)
  \item Evaluation of recurrent VVC (4 or more episodes in 12 months) to confirm the diagnosis by culture or NAAT, identify the \textit{Candida} species, and guide antifungal therapy (fluconazole is preferred for \textit{C. albicans}; \textit{C. glabrata} is inherently less susceptible or resistant to fluconazole and requires topical or alternative therapy)
  \item Investigation of persistent or refractory symptoms despite empiric antifungal therapy, where species identification and antifungal susceptibility testing can guide management
  \item Evaluation of women with risk factors for complicated VVC: uncontrolled diabetes, immunosuppression (HIV, transplant, corticosteroid use), pregnancy, and recent antibiotic use
  \item Evaluation of women with atypical presentations: severe vulvar erythema, edema, excoriations, fissures, or satellite papulopustular lesions (the latter suggestive of candidal folliculitis)
\end{itemize}

\section*{Primary vs Secondary Test}
KOH preparation and saline wet mount are primary bedside diagnostic tests for VVC because they are rapid, simple, and widely available. However, their limited sensitivity (40--70\%) means that a negative KOH/wet mount does NOT exclude VVC. Fungal culture is the primary test when species identification and susceptibility testing are needed (recurrent or refractory cases) and when microscopy is negative with high clinical suspicion. NAAT is the primary diagnostic modality where available, offering near-perfect sensitivity and allowing for simultaneous detection of BV and trichomoniasis.

\section*{How This Test Is Performed}
A sterile, non-lubricated speculum is inserted. The character of the discharge is noted (thick, white, adherent, curd-like). A sterile swab is used to collect a sample of vaginal discharge from the lateral vaginal wall and posterior fornix. The swab is used to: (1) Roll onto a glass slide and a drop of 10\% KOH is added; the slide is examined microscopically under 400x magnification for budding yeast and pseudohyphae. The KOH slide may be gently heated (passed over a flame) to accelerate dissolution of epithelial debris. (2) A second slide with saline for wet mount to evaluate for clue cells, trichomonads, and WBCs. (3) A swab for fungal culture (inoculated onto Sabouraud agar). (4) A swab for NAAT (in the manufacturer-specific transport medium). Vaginal pH should also be measured (normal pH 3.8--4.5 in VVC). Self-collected vaginal swabs have equivalent sensitivity for NAAT. Yeast can also be identified on Gram stain and routine Pap smear, though the sensitivity of Pap smear for diagnosing symptomatic VVC is only approximately 30--50\%; Pap smear identification should be confirmed with culture or NAAT if the patient is symptomatic.

\section*{What Preparation Is Needed}
Patients should avoid douching, using vaginal creams, suppositories, spermicides, or engaging in sexual intercourse for 24--48 hours before specimen collection. Recent use of topical or oral antifungal agents will reduce test sensitivity. Patients should inform the provider if they have applied any vaginal products. The speculum should be inserted with water rather than lubricant, which can interfere with microscopic examination.

\section*{Contraindications and Precautions}
\begin{itemize}
  \item A positive \textit{Candida} culture or NAAT in an asymptomatic woman (no pruritus, burning, or discharge) represents colonization (present in 20--30\% of healthy women) and should NOT be treated. \textit{Candida} species are commensals of the vagina. Treatment of asymptomatic colonization leads to unnecessary antifungal exposure, selection for resistance, and no symptomatic benefit. The presence of pseudohyphae on microscopy is more strongly associated with symptomatic infection than budding yeast alone (pseudohyphae are the invasive, tissue-invasive form), though this distinction is not absolute. \textit{C. glabrata} does not produce pseudohyphae or hyphae and appears only as small budding yeast on microscopy, which makes it easily missed on KOH exam. Approximately 5--10\% of women with VVC have a normal vaginal pH and a negative KOH prep; culture or NAAT should be employed if clinical suspicion remains. Recurrent VVC should prompt evaluation for underlying conditions: diabetes mellitus (HbA1c), HIV, iron deficiency anemia, and immunosuppression. In postmenopausal women not on estrogen therapy, \textit{Candida} vaginitis is extremely rare; atrophic vaginitis is far more likely.
\end{itemize}
\section*{Risks}
Speculum examination and vaginal swab collection are generally well-tolerated. No risk from laboratory tests. The primary clinical risk of VVC is misdiagnosis---treating a non-fungal condition as yeast (leading to continued symptoms and delaying appropriate treatment) or undertreating a complicated infection with a short-course regimen when a longer or alternative regimen is needed.

\section*{What Happens After the Test}
KOH/wet mount results are available immediately (within minutes). Culture: preliminary growth at 24--48 hours; final identification at 48--72 hours. NAAT results in 1--3 hours. Confirmed uncomplicated VVC (mild-to-moderate, sporadic, caused by \textit{C. albicans} in a nonpregnant, immunocompetent host) is treated with a short course of topical azole therapy (clotrimazole cream, miconazole suppository) for 1--7 days, or fluconazole 150 mg PO as a single dose. Complicated VVC (severe, recurrent, non-albicans, pregnancy, uncontrolled diabetes, immunocompromised) requires longer-duration therapy: topical azole for 7--14 days, or fluconazole 150 mg PO every 72 hours for 2--3 doses for severe VVC, or weekly fluconazole 150 mg PO for 6 months for recurrent VVC (maintenance suppressive therapy). For \textit{C. glabrata} VVC, fluconazole is often ineffective; topical boric acid 600 mg intravaginally once daily for 14 days, topical flucytosine 5 g intravaginally at bedtime for 14 days, or topical amphotericin B suppositories may be used. Treatment of sexual partners is NOT recommended (VVC is not sexually transmitted).

\section*{Understanding Results}

\subsection*{Below Normal}
\begin{itemize}
  \item \textbf{KOH/wet mount negative (no yeast, pseudohyphae, or hyphae seen):} Does NOT exclude VVC. Microscopy has sensitivity of only 40--70\%. If clinical suspicion is high (thick white discharge, pruritus, vulvar erythema, normal pH), fungal culture or NAAT should be performed
  \item \textbf{Fungal culture negative:} No \textit{Candida} species isolated. Makes VVC unlikely (culture sensitivity 70--90\%). Consider alternative diagnoses: bacterial vaginosis, trichomoniasis, desquamative inflammatory vaginitis, atrophic vaginitis, contact dermatitis, lichen sclerosus, lichen planus, or physiologic discharge
  \item \textbf{NAAT negative for \textit{Candida} species:} High NPV (>99\%). Essentially excludes candidal vaginitis
\end{itemize}

\subsection*{Above Normal}
\begin{itemize}
  \item \textbf{KOH prep positive (budding yeast with pseudohyphae/hyphae):} Confirms VVC in a symptomatic patient. The presence of pseudohyphae correlates with tissue invasion and symptomatic infection. Treatment should follow. The majority of VVC in otherwise healthy, nonpregnant women is uncomplicated and responds readily to short-course therapy
  \item \textbf{Fungal culture positive for \textit{Candida albicans}:} The most common species (80--92\%). Generally azole-susceptible. If the patient is symptomatic, treat with standard-dose fluconazole or topical azole. If asymptomatic, this represents colonization---do not treat
  \item \textbf{Fungal culture positive for \textit{Candida glabrata}:} The second most common species (5--15\%). Inherently less susceptible to fluconazole; requires higher doses, longer-duration topical therapy, or alternative agents (boric acid, flucytosine). Recurrent VVC caused by \textit{C. glabrata} is particularly challenging
  \item \textbf{Fungal culture positive for \textit{Candida krusei}:} Inherently resistant to fluconazole. Treatment with topical azole, boric acid, or amphotericin B. Rare (<1\% of VVC)
  \item \textbf{NAAT positive for \textit{Candida} species:} Diagnostic of candidal vaginitis in a symptomatic patient. NAAT provides species-level identification (\textit{C. albicans} vs. \textit{C. glabrata} vs. \textit{Candida} species group), guiding antifungal selection. As with culture, do not treat asymptomatic patients with a positive NAAT
  \item \textbf{Vaginal pH in the normal range (3.8--4.5):} Typical of VVC and normal flora. A pH >4.5 in a patient with a positive yeast test suggests co-infection (BV or trichomoniasis) or mixed vaginitis
\end{itemize}

\section*{Normal Results}
KOH/wet mount: \textbf{No yeast, pseudohyphae, or hyphae seen}. Fungal culture: \textbf{No \textit{Candida} species isolated}. NAAT: \textbf{\textit{Candida} species not detected}. Vaginal pH: \textbf{3.8--4.5}.

\section*{Follow-up Tests}
\begin{itemize}
  \item \textbf{Fungal culture with species identification:} For cases where NAAT is not available, KOH is negative but clinical suspicion is high, or recurrent/refractory infection where non-albicans species are suspected
  \item \textbf{Antifungal susceptibility testing (AST):} For \textit{Candida} isolates from patients with recurrent or refractory VVC when fluconazole resistance is suspected. CLSI M27 broth microdilution determines fluconazole MIC. Susceptibility testing can be performed at reference laboratories
  \item \textbf{Multiplex vaginitis panel (NAAT):} Simultaneous testing for BV (\textit{Lactobacillus} depletion, BV-associated bacteria), \textit{Candida} species (including \textit{C. glabrata}), and \textit{T. vaginalis} is ideal for the comprehensive evaluation of vaginitis symptoms
  \item \textbf{HbA1c (glycated hemoglobin):} For women with recurrent VVC to screen for undiagnosed diabetes mellitus, which predisposes to \textit{Candida} overgrowth through elevated vaginal epithelial glycogen and impaired neutrophil function
  \item \textbf{HIV testing:} Recurrent or severe VVC may be a presenting manifestation of HIV infection, particularly in women with other risk factors
  \item \textbf{Vaginal pH and wet mount plus trichomoniasis NAAT:} To exclude concurrent BV or trichomoniasis when VVC is diagnosed but symptoms persist after antifungal therapy
\end{itemize}

\section*{References}
\begin{enumerate}
  \item Workowski KA, Bachmann LH, Chan PA, et al. Sexually transmitted infections treatment guidelines, 2021. MMWR Recomm Rep. 2021;70(4):1-187.
  \item Pappas PG, Kauffman CA, Andes DR, et al. Clinical practice guideline for the management of candidiasis: 2016 update by the IDSA. Clin Infect Dis. 2016;62(4):e1-e50.
  \item Sobel JD. Vulvovaginal candidosis. Lancet. 2007;369(9577):1961-1971.
  \item Nyirjesy P, Brookhart C, Lazenby G, et al. Vulvovaginal candidiasis: a review of the evidence for the 2021 CDC STI Treatment Guidelines. Clin Infect Dis. 2022;74(Suppl 2):S162-S168.
  \item Schwebke JR, Gaydos CA, Nyirjesy P, et al. Diagnostic performance of a molecular test versus clinician assessment of vaginitis. J Clin Microbiol. 2018;56(6):e00252-18.
\end{enumerate}

\clearpage
\section*{Regulatory Status and Authorization Date}
This chapter describes a test category, not a specific commercial product. FDA, CDSCO, EU IVDR, and MHRA approval or registration dates therefore cannot be assigned without the assay manufacturer, product name, instrument, and intended use. For laboratory-developed tests, an individual agency approval date may not exist. See the Preface for the interpretation of regulatory status.

\clearpage
